\documentclass[11pt]{article}
\usepackage[scale=0.7]{geometry}
\usepackage{hyperref}
\usepackage{natbib}
\bibliographystyle{rusnat}

\title{WASP Software Engineering 2026}
\author{Lukas Borggren}
\date{July 2026}

\begin{document}

\maketitle

\section{Introduction}
\label{sec:introduction}
My research is about how large language models (LLMs) can be adapted to and applied in Swedish news media. I am an industrial PhD student at Bonnier News where I have access to – academically speaking – relatively untapped domain expertise and data, as well as real-world use cases. Previously, I have worked as a data scientist at the company, thereby gaining insight into the software engineering (SE) and design challenges in operationalising machine learning (ML) models. My latest research project investigated adaptation of foundational LLMs to Swedish journalism via continued pre-training, using millions of news articles and a domain-specific evaluation benchmark. Currently, I am working on specialising the resulting model(s) for specific editorial tasks, like headline generation and stylistic correction. Arguably the most challenging part of my work is evaluation. Specifically, to what extent to trust automated evaluation, as well as when and how to use human evaluation. Because there are commercial applications of my research output, I am heavily incentivised to design evaluations that translate to actual usability, not just optimise for general metrics and benchmarks. Similarly, I typically know where and how my models will be integrated into existing systems, which informs my research methods in the first place. It also encourages me to involve reporters and editors in the process, as they are generally the end users. For instance, a journalist might desire five diverse article summaries as suggestions, not just a single one; or, our hundreds of newsrooms might each have unique requirements on the formatting and tonality of, say, lead paragraphs. It is not uncommon that these features intertwine research and engineering challenges. For example, it might be monetarily unfeasible to train and serve one LLM per newsroom for some task, and as a result, I need to investigate how newsrooms characteristics can be accounted for already when designing the research method.

\section{Lecture Principles}
\textbf{Behavioural AI Engineering} The topic on how the use of AI tools impacts developers piqued my interest. Superficially because the underlying models are typically LLMs, which is my research area, but fundamentally because it exemplifies a broader phenomenon: the cognitive and social effects of long-term AI tool usage. My own usage of AI tools goes beyond SE, extending to activities like information gathering, brainstorming, and writing. As a result, my professional behaviour has already changed. For example, I can now debug network issues that would have taken me hours in mere minutes, or I can get quick answers to questions by asking a chatbot instead of a senior colleague. Naturally, this almost epistemological shift not only impacts me as an individual, but also has effects on the group level. As more individuals and professions discover how AI tools can aid them, it will inevitably change how teams, departments, and entire organisations are operated. In this broader context, I frequently ponder over how my research is intended to be operationalised. Even though I might not always code the actual implementation, I still believe it is my responsibility to reflect on how the behaviour of end users could be affected. Typically, my research output would power some human-in-the-loop tool that aims to increase the quality and/or efficiency of some aspect of journalists' work. However, if such a tool is sufficiently useful, there is a risk that journalists grow over-reliant on it and stop reviewing its outputs, which could ultimately harm the quality of their work. Conversely, if a tool is suboptimal, the journalists' efficiency might decrease as it becomes more time-consuming to quality check the outputs than to just perform the task manually.

\textbf{Agent Organisation} To me, another interesting topic is how multi-agent systems may be built to mimic human organisation. \citet{heLLMBasedMultiAgentSystems2025} present a path where established SE practices are used to organise agents, from market analysis and role definitions, to industry principles and privacy compliance. This relates to the concept of the "expanding stack", illustrating how the applicability of SE agents has begun to grow beyond code- and prompt-level work to encompass more abstract and managerial tasks. However, it is not obvious to me that recycling existing frameworks is the best way to implement this expanding stack. It is understandable that initial work relies on established theory in areas like SE, organisational theory, and psychology – a structuring that is familiar and transparent to us. However, this is not likely to be the optimal solution. After all, the decades of work underpinning such theory has studied humans, and are ultimately simplified models of how we (should) work and collaborate; by design, they abstract away ubiquitous human traits not characteristic of the specific model, but that might be crucial for its validity. Therefore, to me, it seems overly simplistic to naively repurpose these models for LLM-based agents. Although my own research is not concerned with SE agents, my LLMs are sometimes exposed as specialised components that journalists can include in their own agentic workflows. Commonly, journalists try to model newsroom functions that they are familiar with, configuring complex editorial tasks and roles. Using this approach rarely works as intended, and journalists have to resort to workflow decompositions that are unnatural to them, but that better cater to the strengths of an LLM. In my view, this draws a parallel to designing agent systems for SE, and analogously, maybe one should not stop at the first right answer.

\section{Guest Lecture Principles}
\textbf{SE versus ML Systems} In Julian Frattini's presentation, it is conveyed that traditional requirements engineering techniques in SE may be difficult to apply to ML systems. This echoes the sentiment of \citet{sculleyHiddenTechnicalDebt2015}, that ML systems are distinctly unique compared to other software systems. Personally, I have limited experience of requirements engineering specifically, but I recognise the challenges of applying general SE practices to ML production systems. For example, at my workplace, we have technical principles that guide engineers to design and build systems according to stipulated best practices. The principles are fairly abstract, but even at this level, they are not entirely suitable for ML systems. This can pertain to dependencies between code, models, and data, or characteristics of GPU-accelerated hardware. Granted, this framework is originally not intended to include ML-based solutions, but it is nonetheless notable that they are not automagically covered. To me, this signals that ML systems are not easily reduced to a neat subset of software systems in general. A further challenge is how and when to apply general SE practices to ML research, which distinguishes itself in additional ways. Naturally, research is exploratory and might not have productionisation as an end goal, but SE practices can still facilitate aspects like collaboration, security, efficiency, and reproducibility – all important for the quality of research.

\textbf{Persistent Change} Per Lenberg raised an interesting implication of the current redefinition of the SE role: it will persistently change and developers have to get used to it. To me, this highlights an interesting demand for a new type of adaptability. Software engineers are already used to change – new tools to use, new principles for arranging code, new architectures for infrastructure setups. But currently, the change is typically limited to a technical silo, and importantly, the silos themselves do not change. I believe the current development points to a more existential dynamism, that is, the entire role definition transforms, not just aspects of it. A profound effect of this might be a change in engineers' self-image. Up until today, the role has been highly associated with technical expertise, with deep knowledge and care about some parts of the tech stack. However, this view is challenged as LLM agents are increasingly able to replace parts of this expertise, forcing developers to evolve their skill portfolio. From being capable to perform junior developer tasks, agents are now becoming competitive with senior experts. This impacts both students, that have yet to start working, and experienced engineers that have coded for decades. What does the future hold in store for them? Moreover, I do not believe this phenomenon is limited to SE. In my vicinity,  I see roles like investigate journalist, data scientist, and researcher changing rapidly. Frankly, I can only guess how my own work as an NLP researcher will look like in a decade.

\section{Data Scientists, Software Engineers \& AI Engineers}
The book generally describes data scientists as focused on exploratory model training and data curation, and less interested in code quality or deployment. In contrast, software engineers are described as focused on delivering according to requirements, working on aspects like usability, security, and costs. Personally, I definitely recognise these tropes and think there are some merits to them. However, in reality, I believe the boundary between the roles is fuzzy and conditioned on several factors. Firstly, there are role semantics. A large company might have dedicated research engineers, data scientist, ML engineers etcetera, whereas a smaller tech organisation might have one generalist data-science role covering the entire stack. Having worked as data scientist in the latter context, I know how small of a fraction model tinkering can be of the total work, which also encompasses tasks like data engineering, testing, deployment, and monitoring. Secondly, agreeing with the book, I believe that even highly specialised data scientists and researchers have to adopt a somewhat systemic view of development. For instance, requirements on latency or scaling might directly interweave operational costs and model development. In my own research, I have had to investigate various model-compression techniques to meet latency requirements. Ensuring that services can auto-scale down to zero replicas, I have considered cold-start problems related to model architecture and loading. Another example is that from Section \ref{sec:introduction}: it might be unfeasible to host hundreds of specialised models, and I have to figure out an alternative training method to obtain competitive results using a single model.

Increasingly, I think the two roles are merging. Each will continue having specialities, but as ML becomes ever-more mature and productionised, it will inevitably be more tightly integrated with SE and the roles will have to learn from each other. I also believe the access to capable commercial and open-weight foundational models boosts this developments; both because it makes ML part of even more systems, but also because developers use them as tools in their own work. After all, integrating towards existing LLMs requires virtually no data science expertise, but it might be good if the developer has some idea about the possibilities and limitations of the technology. Similarly, to be efficient in their own work, software engineers are encouraged to learn more about ML just to be able to use coding agents more proficiently. As a result, model training has become an even smaller piece of ML products, while the surrounding engineering work remains and evolves. After all, today, it might be more valuable to business that employees can prompt and harness existing models, than that they can train new models. As to naming these emerging cross-pollination of expertises, I have no clear idea. I have some preconceived SE4ML notion about what AI/ML engineers do, even though I find the names rather nondescript. For ML4SE on the other hand, I have no firm grasp on roles yet. In the end, with the growing supply of capable models, ML is really becoming another tool in an expanding toolbox. After all, a reporter writing their story on a laptop instead of a typewriter is not consider a computer expert – at least not any more.

\section{Paper Analysis}
\subsection{MLScent: A Tool for Anti-Pattern Detection in ML Projects}
\label{sec:mlscent}
\citet{shivashankarMLScentToolAntiPattern2025} propose MLScent, a novel analysis tool for detecting ML-specific anti-patterns and code smells. Fundamentally, the authors motivate their work analogously to \citet{sculleyHiddenTechnicalDebt2015}; namely, that compared to SE in general, ML systems present unique challenges and thus requires equally unique tooling. Operating on the code level, MLScent implements both general and framework-specific detectors unique to ML and data science, targeting anti-patterns and code smells compiled from prior work. This focus and comprehensiveness make the tool novel in contrast to existing solutions. The tool is based on static analysis with abstract syntax trees and is implemented as a CLI, outputting both human- and machine-readable formats. Comparing MLScent's output to that of human annotators, there is strong agreement, and qualitative feedback generally suggests that the tool is useful.

The research behind MLScent is not that related to my own area, but I believe the tool itself could be useful to me. Particularly so if integrated into existing development platforms, and the authors state that they indeed plan to develop IDE plugins for VSCode, and CI/CD integrations. Writing experiment code for my own research, I would happily use these implementations to increase code quality and capture potential issues at an early stage. It would be a natural extension to general SE tooling that I already use for Python, like Pylint and MyPy. This could be beneficial from a pure efficiency standpoint, such as decreased memory usage, but also for avoiding bugs or unintentional data modifications that might affect the actual research results.

Rather than describing a hypothetical AI-intensive software project, I believe the integration and usefulness of MLScent can be highlighted by virtually any ML codebase. Specifically, I believe the benefits will increase as the scale of such projects grows, both in terms of the size and complexity of the code itself, but also the number of contributors and project matureness. For a production-level monorepo, MLScent could act as a valuable additional step of verification in CI/CD pipelines, maintaining coding standards and best-practices specifically for the ML-centric code. Not only would this enforce uniform standards, but also explicitly show junior developers how code is expected to be written. The creators of MLScent further claim the tool to be extensible, meaning it provides developers with a framework for how to implement their own detectors. Downstream, the usability of the tool could manifest itself as less time spent on debugging erroneous deployments, and reduced on-demand infrastructure costs if code runs more efficiently.

\subsection{Machine Learning Observability in Practice}
\citet{leestMonitoringObservabilityMachine2025} present and systemise empirical results on how ML systems are monitored in practice, and identify gaps in current observability practices. The authors motivate their study by noting the unique challenges of monitoring ML systems, with existing tooling being compartmentalised and frameworks being purely theoretically derived. In contrast, the paper seeks to find empirical evidence for what information ML practitioners use for monitoring, and to what extent it is systemised. Via focus group sessions with practitioners, answers are collected and analysed using existing frameworks. Subsequently, the authors group and highlight common practices, both conceptually and concretely, additionally pointing to non-mature monitoring areas. The emerging image showcases the system-level intricacies of ML applications, from data sources to business metrics, and the difficulties of implementing monitoring, as well as debugging.

Again, this research is not directly connected to my own, but I recognise many of the common themes identified for production-level ML systems. Because the output is not as concrete as MLScent, Section \ref{sec:mlscent}, the application of this paper's ideas is more implicit, perhaps even a mindset. I believe that some of these principles are useful even when writing research code. Even though many activities in research are one-offs for a particular paper, observability is still useful for analysis and debugging. For instance, I implement data preprocessing in production-levels pipelines, decomposing actions into atomic steps and saving intermediate outputs to enable traceability. This is helpful to both structure the code and simplify debugging, for example, if a lot of data is dropped – I can trace it down to a specific pipeline step. For me, model training is more monolithic because I am confined to academic infrastructure for that step. Either way, I find these concepts useful for any research that will be applied. It makes me design and structure solutions with our production systems in mind, lowering the threshold for any downstream productionalisation.

Similar to MLScent, I believe the paper's ideas could be useful in virtually any AI-intensive SE project. The authors provide a framework to sanity check implementations against, essentially allowing benchmarking of observability practices against industry standards. The challenge would be how to operationalise the comparison. The authors give some examples on tools and frameworks used for specific monitoring areas, but it may be difficult to determine which what suits a specific project. Furthermore, the challenging use-cases are the ones where even seasoned practitioners do not have systemised approaches. However, such challenges, like gaps between system parts, could be informative already in designing the system, that is, trying to minimise these error-prone areas in the first place.

\section{Evidence, Originality \& GenAI-Use Transparency}
All references except the course materials used are listed below. They are all peer-reviewed papers, although only the preprint version of \citet{leestMonitoringObservabilityMachine2025} is available at the moment. As a result, I deem the references trustworthy. I used ChatGPT and Gemini to search for relevant papers and to polish text.

\bibliography{references}

\end{document}
